\documentclass[UTF8,a4paper,11pt]{ctexart}
\usepackage[margin=2.0cm]{geometry}
\usepackage{amsmath,amssymb,amsfonts,mathtools,bm}
\usepackage{xcolor}
\usepackage{enumitem}
\usepackage{hyperref}
\hypersetup{colorlinks=true,linkcolor=blue,urlcolor=blue}

% 标注命令
\newcommand{\bp}[1]{\textcolor{red}{\textbf{【要背】}}\ #1}
\newcommand{\uf}[1]{\textcolor{blue}{\textbf{【要了解】}}\ #1}
\newcommand{\imp}[1]{\textcolor{teal}{\textbf{（考点）}}\ #1}

\title{\textbf{理论力学复习提纲（详细版）}}
\author{燕山大学 · 应用物理学 · 理论力学\\ 据课堂讲义、手写笔记及往届考题整理}
\date{}

\begin{document}
\maketitle

\section*{0. 回忆考题 — 考点对应表}
\begin{center}
\begin{tabular}{|l|l|l|}
\hline
\textbf{题型} & \textbf{回忆到的题目} & \textbf{对应章节} \\
\hline
概念题1 & 完整约束还是非完整约束 & 第1节 约束分类 \\
概念题2 & 相空间 & 第4节 相空间 \\
概念题3 & 简正坐标 & 第3节 小振动 \\
概念题4 & 虚功原理 & 第1节 虚功原理 \\
概念题5 & 简单概念（未想起） & 全章基础概念 \\
\hline
选择1 & 完整约束/非完整约束 & 第1节 \\
选择2 & 拉氏方程缺 $\dot q$ 项 $\rightarrow$ 选广义动量守恒 & 第2节 \\
选择3 & 哈密顿原理变量 $q,p,t$ & 第5节 \\
选择4、5 & 简单（未想起） & — \\
\hline
大题1 & 耦合摆证明 & 第3节 \\
大题2 & 母函数证明 & 第5节 \\
大题3 & 滑轮题：传统力学/拉氏/哈密顿分别解 & 第2、4节 \\
大题4 & 两自由度拉氏函数（滑轮，计算难） & 第2、4节 \\
大题5 & 虚功原理课本例题 & 第1节 \\
\hline
\end{tabular}
\end{center}

\tableofcontents
\newpage

% ===================== 第1节 =====================
\section{约束与虚功原理}
\imp{概念题1、4，选择题1，大题5均出自本节，是整门课的入口。}

\subsection{约束的分类 \bp{}}
设系统由 $N$ 个质点组成，受 $k$ 个约束，自由度为 $s=3N-k$（完整约束时）。
\begin{enumerate}[itemsep=2pt]
  \item \bp{几何约束（位置约束）}：只含坐标与时间，$f(\bm r_i,t)=0$。
  \item \uf{运动学约束（速度约束）}：含坐标与速度，$\varphi(\bm r_i,\dot{\bm r}_i,t)=0$。若能积分成有限形式则退化为几何约束。
  \item \bp{定常约束}：约束方程不显含时间 $t$；\bp{非定常约束}：显含 $t$。
  \item \bp{完整约束}：可写成坐标有限形式的约束（含可积分的速度约束）。系统位形可由 $s$ 个广义坐标完全确定。\\
        \bp{非完整约束}：不可积分的速度约束（如纯滚动、冰刀侧滑）。此时自由度 $\neq$ 广义坐标数。
  \item \uf{单侧约束}：不等式 $f\ge 0$（如球不离开桌面）；\uf{双侧约束}：等式 $f=0$（如刚性杆长固定）。
\end{enumerate}
\bp{判据（必考）}：看约束能否写成仅关于坐标（不含量纲为速度的微分关系）的有限方程。能量守恒式、运动学速度式若不可积 $\Rightarrow$ 非完整。

\subsection{虚位移 \bp{}}
在\textbf{给定瞬时}、令 $\delta t=0$（时间“冻结”），约束所允许的、无限小的位置变更 $\delta\bm r_i$。它满足约束方程的线性化：
\[
\sum_i \frac{\partial f}{\partial \bm r_i}\cdot \delta\bm r_i=0.
\]
\bp{虚位移与实位移区别}：实位移 $d\bm r_i$ 发生在 $dt$ 内、受动力学支配；虚位移 $\delta\bm r_i$ 是纯几何概念、与 $dt$ 无关。

\subsection{虚功与理想约束 \bp{}}
\begin{itemize}[itemsep=2pt]
  \item 虚功：$\delta W=\sum_i \bm F_i\cdot\delta\bm r_i$。
  \item \bp{理想约束}：约束力 $\bm N_i$ 在\textbf{任意}虚位移上总虚功为零
  \[
  \sum_i \bm N_i\cdot\delta\bm r_i=0.
  \]
  光滑接触、不可伸长绳、刚性杆、固定铰均满足。
\end{itemize}

\subsection{虚功原理 \bp{}}
\bp{平衡的充分必要条件}（理想约束下）：\textbf{所有主动力的虚功之和为零}
\[
\boxed{\sum_i \bm F_i^{(a)}\cdot\delta\bm r_i=0}\qquad\text{或分量式}\qquad
\sum_i (F_{ix}\delta x_i+F_{iy}\delta y_i+F_{iz}\delta z_i)=0.
\]
\uf{理解}：虚功原理把平衡问题化为纯几何/代数问题，无需考虑约束力（理想约束已被消去）。

\subsection{达朗贝尔原理 \uf{}}
引入惯性力 $-\bm m_i\bm a_i$，平衡方程变为
\[
\sum_i (\bm F_i^{(a)}-\bm m_i\bm a_i)\cdot\delta\bm r_i=0.
\]
这是拉格朗日方程的出发点。

\subsection{例题：两刚性杆光滑铰链平衡 \uf{}}
\uf{思路}（课本/笔记例题）：设两等长杆以铰链连于墙，在重力与端点外力下平衡。
\begin{enumerate}[itemsep=2pt]
  \item 取广义坐标（如各杆与竖直夹角 $\theta_1,\theta_2$），写出各主动力（重力、外力）作用点的虚位移；
  \item 由虚位移关系 $\delta\bm r=f(\delta\theta_1,\delta\theta_2)$ 代入 $\sum\bm F^{(a)}\cdot\delta\bm r=0$；
  \item 因 $\delta\theta_1,\delta\theta_2$ 独立，各自系数为零 $\Rightarrow$ 平衡方程。
\end{enumerate}
\imp{大题5即此类，按“选坐标→写虚位移→代虚功原理→系数归零”四步走即可。}

% ===================== 第2节 =====================
\section{拉格朗日方程}
\imp{选择题2（缺 $\dot q$ 项选广义动量）、大题3/4（滑轮）出自本节。}

\subsection{广义坐标与自由度 \bp{}}
\bp{广义坐标} $q_1,\dots,q_s$：完全确定系统位形的最少独立参数，数目 $=s$（完整约束）。
\bp{自由度}：确定系统位形所需的独立坐标数。

\subsection{拉格朗日函数 \bp{}}
\[
\boxed{L(q,\dot q,t)=T-V}
\]
其中 $T$ 为动能，$V$ 为势能。保守力下 $V$ 不含广义速度。

\subsection{拉格朗日方程（第二类） \bp{}}
对每一广义坐标 $q_j\ (j=1,\dots,s)$：
\[
\boxed{\frac{d}{dt}\!\left(\frac{\partial L}{\partial \dot q_j}\right)-\frac{\partial L}{\partial q_j}=Q_j}
\]
保守力时 $Q_j=0$，化为
\[
\boxed{\frac{d}{dt}\!\left(\frac{\partial L}{\partial \dot q_j}\right)-\frac{\partial L}{\partial q_j}=0.}
\]
\bp{解题四步}：（1）选广义坐标；（2）写 $T,V$ 得 $L$；（3）代入拉氏方程；（4）化简求运动方程。

\subsection{广义动量与循环坐标 \bp{}}
\[
p_j=\frac{\partial L}{\partial \dot q_j}\quad\text{（广义动量）}.
\]
\bp{循环坐标}：若 $L$ 不含某 $q_k$（$\partial L/\partial q_k=0$），则对应 $p_k=\text{const}$（广义动量守恒）。这就是选择题2的考点：拉氏方程缺 $\dot q$ 项（实为缺 $q$ 本身）时，对应广义动量守恒。

\subsection{广义能量积分 \uf{}}
若约束定常且 $L$ 不显含 $t$，则
\[
h=\sum_j \dot q_j\frac{\partial L}{\partial \dot q_j}-L=T+V=E=\text{const}.
\]
\uf{理解}：雅可比积分 = 机械能守恒（保守系统、定常约束）。

\subsection{由达朗贝尔推拉氏方程 \uf{}}
\uf{关键三步}：（1） $\delta\bm r_i=\sum_j\frac{\partial\bm r_i}{\partial q_j}\delta q_j$；（2） 利用 $\frac{\partial\dot{\bm r}_i}{\partial\dot q_j}=\frac{\partial\bm r_i}{\partial q_j}$ 与 $\frac{d}{dt}\frac{\partial\bm r_i}{\partial q_j}=\frac{\partial\dot{\bm r}_i}{\partial q_j}$；（3） 得 $\frac{d}{dt}\frac{\partial T}{\partial\dot q_j}-\frac{\partial T}{\partial q_j}=Q_j$，保守力下换成 $L$。

\subsection{例题：阿特伍德机（拉格朗日法） \uf{}}
质量 $M>m$ 跨过无质量滑轮，绳长固定。取 $x$ 为 $M$ 下降距离：
\[
T=\tfrac12(M+m)\dot x^2,\qquad V=-(M-m)gx\quad(\text{取下降为正}).
\]
\[
L=\tfrac12(M+m)\dot x^2+(M-m)gx.
\]
代入拉氏方程：
\[
\frac{d}{dt}\big((M+m)\dot x\big)-(M-m)g=0
\;\Rightarrow\;
\boxed{\ddot x=\frac{M-m}{M+m}g.}
\]

\subsection{例题：耦合摆（见第3节）}

% ===================== 第3节 =====================
\section{小振动与简正坐标}
\imp{概念题3（简正坐标）、大题1（耦合摆证明）出自本节。}

\subsection{多自由度小振动 \uf{}}
平衡位形附近，$T=\tfrac12\sum a_{jk}\dot q_j\dot q_k$，$V=\tfrac12\sum b_{jk}q_jq_k$（二次型）。引入简正坐标前，运动方程为耦合的线性微分方程组。

\subsection{简正坐标 \bp{}}
\bp{定义}：通过线性变换 $\bm\xi=\bm A\bm q$，使动能、势能同时对角化：
\[
T=\tfrac12\sum_k \dot\xi_k^2,\qquad V=\tfrac12\sum_k \omega_k^2\xi_k^2.
\]
对角化后的坐标 $\xi_k$ 即\textbf{简正坐标}，每个只对应一个独立的简谐振动（频率 $\omega_k$）。

\subsection{耦合摆本征频率（证明题） \uf{}}
两相同单摆，长 $l$、质量 $m$，中点弹簧劲度 $k$。小角近似：
\[
T=\tfrac12ml^2(\dot\theta_1^2+\dot\theta_2^2),\quad
V=\tfrac12mgl(\theta_1^2+\theta_2^2)+\tfrac12kl^2(\theta_1-\theta_2)^2.
\]
拉氏方程给出
\[
ml^2\ddot\theta_1=-mgl\,\theta_1-kl^2(\theta_1-\theta_2),\qquad
ml^2\ddot\theta_2=-mgl\,\theta_2-kl^2(\theta_2-\theta_1).
\]
\bp{取简正坐标} $\xi_1=\theta_1+\theta_2,\ \xi_2=\theta_1-\theta_2$，代入得
\[
\ddot\xi_1=-\frac{g}{l}\,\xi_1,\qquad
\ddot\xi_2=-\left(\frac{g}{l}+\frac{2k}{m}\right)\xi_2.
\]
\bp{两个简正模频率}：
\[
\boxed{\omega_1=\sqrt{\frac{g}{l}}}\quad\text{（对称模，同相）},\qquad
\boxed{\omega_2=\sqrt{\frac{g}{l}+\frac{2k}{m}}}\quad\text{（反对称模，反相）}.
\]
\imp{大题1即此证明：核心是“选 $\xi_1=\theta_1+\theta_2,\ \xi_2=\theta_1-\theta_2$ 解耦”。}

% ===================== 第4节 =====================
\section{哈密顿力学}
\imp{概念题2（相空间）、选择题3（哈密顿原理变量）、大题3（哈密顿解滑轮）出自本节。}

\subsection{勒让德变换与哈密顿量 \bp{}}
由拉氏量作勒让德变换，引入广义动量 $p_j=\partial L/\partial\dot q_j$，定义
\[
\boxed{H(q,p,t)=\sum_j p_j\dot q_j-L.}
\]
\bp{物理意义}：保守、定常系统下 $H=T+V=E$。

\subsection{正则方程 \bp{}}
\[
\boxed{\dot q_j=\frac{\partial H}{\partial p_j},\qquad
\dot p_j=-\frac{\partial H}{\partial q_j}}\qquad(j=1,\dots,s).
\]
共 $2s$ 个一阶方程，是哈密顿形式的核心。

\subsection{相空间 \bp{}}
\bp{定义}：以 $s$ 个广义坐标 $q_j$ 与 $s$ 个广义动量 $p_j$ 为坐标轴构成的 $2s$ 维空间。系统的运动状态由相空间中的一个\textbf{相点} $(q,p)$ 表示；运动对应一条相轨迹。

\subsection{泊松括号 \bp{}}
对任意两个力学量 $u,v$：
\[
\boxed{\{u,v\}=\sum_j\left(\frac{\partial u}{\partial q_j}\frac{\partial v}{\partial p_j}
-\frac{\partial u}{\partial p_j}\frac{\partial v}{\partial q_j}\right)}.
\]
\bp{基本性质（必记）}：
\begin{enumerate}[itemsep=1pt]
  \item 反对称：$\{u,v\}=-\{v,u\}$；
  \item 双线性；
  \item 乘积法则：$\{uv,w\}=u\{v,w\}+\{u,w\}v$；
  \item \bp{雅可比恒等式}：$\{\{u,v\},w\}+\{\{v,w\},u\}+\{\{w,u\},v\}=0$；
  \item \bp{正则括号}：$\{q_j,q_k\}=0,\ \{p_j,p_k\}=0,\ \{q_j,p_k\}=\delta_{jk}$；
  \item 运动方程：$\dot u=\{u,H\}+\dfrac{\partial u}{\partial t}$。
\end{enumerate}

\subsection{刘维尔定理 \uf{}}
相空间中代表点密度 $\rho(q,p,t)$ 满足
\[
\frac{\partial\rho}{\partial t}+\sum_j\left(\frac{\partial(\rho\dot q_j)}{\partial q_j}
+\frac{\partial(\rho\dot p_j)}{\partial p_j}\right)=0
\;\Rightarrow\;\frac{d\rho}{dt}=0.
\]
\uf{含义}：相体积在演化中守恒（不可压缩流）。

\subsection{位力定理 \uf{}}
对束缚系统（时间平均 $\langle\cdot\rangle$）：
\[
2\langle T\rangle=\left\langle\sum_i \bm F_i\cdot\bm r_i\right\rangle.
\]
若 $V\propto r^n$ 则 $2\langle T\rangle=n\langle V\rangle$。

\subsection{例题：阿特伍德机（哈密顿法） \uf{}}
同第2节：$p=\dfrac{\partial L}{\partial\dot x}=(M+m)\dot x$，
\[
H=\frac{p^2}{2(M+m)}-(M-m)gx.
\]
正则方程：
\[
\dot x=\frac{\partial H}{\partial p}=\frac{p}{M+m},\qquad
\dot p=-\frac{\partial H}{\partial x}=(M-m)g.
\]
两式消去 $p$ 得 $\ddot x=\dfrac{M-m}{M+m}g$，与拉氏法、牛顿法一致。
\imp{大题3要求三种方法对照，记住“牛顿法看受力、拉氏法看 $L$、哈密顿法看 $H$ 与正则方程”即可。}

% ===================== 第5节 =====================
\section{变分原理与正则变换}
\imp{选择题3（哈密顿原理变量 $q,p,t$）、大题2（母函数证明）出自本节。}

\subsection{泛函与欧拉—拉格朗日方程 \uf{}}
变分问题 $\delta\!\int_{x_1}^{x_2}F(x,y,y')\,dx=0$ 给出
\[
\frac{\partial F}{\partial y}-\frac{d}{dx}\frac{\partial F}{\partial y'}=0.
\]

\subsection{哈密顿原理 \bp{}}
\bp{内容}：在固定端点 $\delta q(t_1)=\delta q(t_2)=0$ 的真实运动轨道上，作用量
\[
S=\int_{t_1}^{t_2}L(q,\dot q,t)\,dt
\]
取\textbf{驻值}：
\[
\boxed{\delta S=\delta\int_{t_1}^{t_2}L\,dt=0.}
\]
\bp{变量}：作用量泛函依赖于 $q(t),\dot q(t),t$；哈密顿原理是导出拉氏方程、正则方程的变分基础。选择题3考的“变量 $q,p,t$”实际指哈密顿正则方程/相空间的变量，注意区分“哈密顿原理的变分变量 $q,\dot q,t$”与“正则方程的变量 $q,p,t$”。

\subsection{正则变换 \uf{}}
从 $(q,p)$ 到 $(Q,P)$ 的变换，若存在新哈密顿量 $K$ 使
\[
\dot Q_j=\frac{\partial K}{\partial P_j},\qquad
\dot P_j=-\frac{\partial K}{\partial Q_j},
\]
则为正则变换。它保持力学形式不变，是求解复杂系统的利器。

\subsection{四类母函数 \bp{}}
设新哈密顿量 $K=H+\dfrac{\partial F}{\partial t}$，四类生成函数（母函数）及变换关系：
\begin{enumerate}[itemsep=2pt]
  \item \bp{第一类 $F_1(q,Q,t)$}：
  \[
  p_j=\frac{\partial F_1}{\partial q_j},\quad P_j=-\frac{\partial F_1}{\partial Q_j}.
  \]
  \item \bp{第二类 $F_2(q,P,t)=F_1+\sum Q_jP_j$}：
  \[
  p_j=\frac{\partial F_2}{\partial q_j},\quad Q_j=\frac{\partial F_2}{\partial P_j}.
  \]
  \item \bp{第三类 $F_3(p,Q,t)=F_1-\sum q_jp_j$}：
  \[
  q_j=-\frac{\partial F_3}{\partial p_j},\quad P_j=-\frac{\partial F_3}{\partial Q_j}.
  \]
  \item \bp{第四类 $F_4(p,P,t)=F_1-\sum q_jp_j+\sum Q_jP_j$}：
  \[
  q_j=-\frac{\partial F_4}{\partial p_j},\quad Q_j=\frac{\partial F_4}{\partial P_j}.
  \]
\end{enumerate}
\bp{记忆法（只需背第3步）}：
\begin{enumerate}[itemsep=2pt]
  \item \textbf{自变量}：四类即 $(q,Q),(q,P),(p,Q),(p,P)$ 四种组合——
  $F_1$ 留旧坐标$q$新坐标$Q$，$F_2$ 留旧$q$新动量$P$，$F_3$ 留旧动量$p$新$Q$，$F_4$ 留旧$p$新$P$。
  \item \textbf{加减项}（从 $F_1$ 做勒让德变换）：把 $Q$ 换成 $P$ 时\textbf{加} $\sum Q_jP_j$；把 $q$ 换成 $p$ 时\textbf{减} $\sum q_jp_j$。故
  $F_2=F_1+\sum Q_jP_j,\ F_3=F_1-\sum q_jp_j,\ F_4=F_1-\sum q_jp_j+\sum Q_jP_j$。
  \item \textbf{偏导符号（最关键、只背这条）}：母函数\textbf{显含哪个变量}，就对谁求偏导得到其共轭量，符号恒固定：
  \[
  \begin{array}{c|c}
  \text{母函数显含} & \text{偏导关系} \\ \hline
  q_j & p_j=+\dfrac{\partial F}{\partial q_j} \\[6pt]
  Q_j & P_j=-\dfrac{\partial F}{\partial Q_j} \\[6pt]
  p_j & q_j=-\dfrac{\partial F}{\partial p_j} \\[6pt]
  P_j & Q_j=+\dfrac{\partial F}{\partial P_j}
  \end{array}
  \]
  即：坐标对 $(q,Q)$ 恒为“$p=+\partial/\partial q,\ P=-\partial/\partial Q$”；动量对 $(p,P)$ 恒为“$q=-\partial/\partial p,\ Q=+\partial/\partial P$”。
\end{enumerate}

\subsection{哈密顿—雅可比方程 \uf{}}
取 $F_2=S(q,P,t)$ 且令新动量 $P_j=\alpha_j=\text{const}$，则新坐标 $Q_j=\beta_j=\text{const}$，母函数 $S$ 满足
\[
\boxed{H\!\left(q,\frac{\partial S}{\partial q},t\right)+\frac{\partial S}{\partial t}=0.}
\]
解出 $S$ 后：
\[
p_j=\frac{\partial S}{\partial q_j},\qquad \beta_j=\frac{\partial S}{\partial\alpha_j}.
\]
\imp{大题2“母函数证明”通常要求：写出某类母函数定义 $\rightarrow$ 证明变换后的方程仍为正则形式（即 $K$ 存在且 $K=H+\partial F/\partial t$），或验证一组给定的 $(Q,P)$ 变换是正则的（检查泊松括号 $\{Q_j,Q_k\}=0,\{P_j,P_k\}=0,\{Q_j,P_k\}=\delta_{jk}$）。}

% ===================== 第6节 =====================
\section{考场策略与易错点}
\begin{enumerate}[itemsep=3pt]
  \item \bp{大题按步给分}：哪怕算不出最终数，也要写出“选坐标 $\rightarrow$ 写 $L$（或 $H$）$\rightarrow$ 代方程 $\rightarrow$ 列式”，步骤分占大头。
  \item \uf{小振动}：先写 $T,V$ 二次型，再令 $\xi=\theta_1\pm\theta_2$ 解耦；计算易错在 $2k/m$ 系数，检查量纲。
  \item \uf{母函数正负号}：$p_j=+\partial F/\partial q_j$，$P_j=-\partial F/\partial Q_j$（第一类）；$F_2$ 用 $+\partial F_2/\partial P_j=Q_j$。正负号是高频失分点。
  \item \bp{广义动量守恒}：拉氏方程“不显含某 $q_k$”（非“不显含 $\dot q_k$”）$\Rightarrow p_k$ 守恒。选择题2即此陷阱。
  \item \uf{三种方法对照}（滑轮题）：牛顿法看受力图；拉氏法 $L=T-V$ 一步出运动方程；哈密顿法先求 $p$ 再写 $H$ 与正则方程。
  \item \bp{完整 vs 非完整}：考概念题/选择题时，先判断“能否写成仅坐标的有限方程”。纯滚动、轮滑冰刀是非完整典型。
\end{enumerate}

\end{document}
